% ---------- lingua e codifica ----------
\usepackage[T1]{fontenc}
\usepackage[utf8]{inputenc}
\usepackage[italian]{babel}
\usepackage{lmodern}
\usepackage{microtype}

% ---------- pagina ----------
\usepackage[a4paper,margin=2.5cm]{geometry}
\usepackage{fancyhdr}
\pagestyle{fancy}
\fancyhf{}
\fancyhead[LE,RO]{\thepage}
\fancyhead[RE]{\small\nouppercase{\leftmark}}
\fancyhead[LO]{\small\nouppercase{\rightmark}}
\setlength{\headheight}{14pt}

% ---------- matematica ----------
\usepackage{amsmath,amssymb,amsthm,mathtools}
\usepackage{bm}
\usepackage[per-mode=symbol]{siunitx}
\sisetup{output-decimal-marker={,}, exponent-product=\cdot, group-separator={\,}}

% ---------- grafica ----------
\usepackage{graphicx}
\usepackage{float}
\usepackage{tikz}
\usetikzlibrary{arrows.meta,calc,positioning,decorations.pathreplacing,patterns,angles,quotes}
\usepackage{pgfplots}
\pgfplotsset{compat=1.18}
\usepgfplotslibrary{fillbetween,groupplots}
\usepackage{subcaption}
\usepackage{booktabs,tabularx,multirow}
\usepackage{enumitem}
\usepackage[font=small,labelfont=bf]{caption}

% ---------- box colorati ----------
\usepackage{etoolbox,needspace}
\usepackage[most]{tcolorbox}
\tcbset{enhanced, breakable, fonttitle=\bfseries, arc=2pt, before skip=11pt plus 2pt, after skip=11pt plus 2pt}
\newcommand{\sottotitolo}[1]{\if\relax\detokenize{#1}\relax\else: #1\fi}
\newtcolorbox{esempio}[1][]{unbreakable, colback=blue!3, colframe=blue!50!black, title={Esempio\sottotitolo{#1}}}
\newtcolorbox{attenzione}[1][]{unbreakable, colback=red!3, colframe=red!60!black, title={Attenzione\sottotitolo{#1}}}
\newtcolorbox{richiamo}[1][]{colback=green!3, colframe=green!40!black, title={Richiamo di matematica\sottotitolo{#1}}}
\newtcolorbox{prerequisiti}{unbreakable, colback=yellow!8, colframe=orange!60!black, title={Prerequisiti (vedi Capitolo~0)}, fontupper=\small}
\newtcolorbox{riquadro}[1]{unbreakable, colback=blue!4, colframe=blue!55!black, title={#1}, fontupper=\normalsize}
\newcommand{\serve}[1]{\par\medskip\noindent\textbf{Perch\'e ci serve in questo corso.} #1\par\medskip}
\newtcolorbox{riassunto}[1][]{unbreakable, colback=gray!5, colframe=gray!60!black, title={In sintesi\sottotitolo{#1}}}

% ---------- codice ----------
\usepackage{listings}
\lstset{language=Matlab, basicstyle=\ttfamily\footnotesize, keywordstyle=\color{blue!70!black},
        commentstyle=\color{green!40!black}, stringstyle=\color{red!60!black},
        numbers=left, numberstyle=\tiny\color{gray}, numbersep=6pt, xleftmargin=2.2em, xrightmargin=1.2em, framexleftmargin=2em, framexrightmargin=2pt, frame=single, breaklines=true, breakindent=0pt, postbreak=\mbox{\textcolor{gray}{$\hookrightarrow$}\space}, aboveskip=10pt, belowskip=10pt, showstringspaces=false,
        literate={è}{{\`e}}1 {à}{{\`a}}1 {ù}{{\`u}}1 {ò}{{\`o}}1 {é}{{\'e}}1}

% ---------- riferimenti ----------
\usepackage[hidelinks]{hyperref}
\usepackage[italian,nameinlink]{cleveref}

% ---------- teoremi ----------
\theoremstyle{definition}
\newtheorem{definizione}{Definizione}[chapter]
\theoremstyle{plain}
\newtheorem{teorema}{Teorema}[chapter]
\newtheorem{proposizione}[teorema]{Proposizione}
\newtheorem{lemma}[teorema]{Lemma}
\newtheorem{proprieta}[teorema]{Propriet\`a}
\theoremstyle{remark}
\newtheorem{osservazione}{Osservazione}[chapter]
\renewcommand{\proofname}{Dimostrazione}
\crefname{teorema}{teorema}{teoremi}           \Crefname{teorema}{Teorema}{Teoremi}
\crefname{proposizione}{proposizione}{proposizioni} \Crefname{proposizione}{Proposizione}{Proposizioni}
\crefname{proprieta}{propriet\`a}{propriet\`a} \Crefname{proprieta}{Propriet\`a}{Propriet\`a}
\crefname{lemma}{lemma}{lemmi}                 \Crefname{lemma}{Lemma}{Lemmi}
\crefname{definizione}{definizione}{definizioni} \Crefname{definizione}{Definizione}{Definizioni}
\crefname{osservazione}{osservazione}{osservazioni} \Crefname{osservazione}{Osservazione}{Osservazioni}
\crefname{appendix}{appendice}{appendici} \Crefname{appendix}{Appendice}{Appendici}

% ---------- riquadri per definizioni e teoremi ----------
\tcolorboxenvironment{definizione}{unbreakable, colback=blue!3, colframe=blue!55!black, boxrule=0.6pt, left=6pt, right=6pt, top=3pt, bottom=3pt, before skip=11pt plus 2pt, after skip=11pt plus 2pt}
\tcolorboxenvironment{teorema}{unbreakable, colback=orange!5, colframe=orange!70!black, boxrule=0.6pt, left=6pt, right=6pt, top=3pt, bottom=3pt, before skip=11pt plus 2pt, after skip=11pt plus 2pt}
\tcolorboxenvironment{proposizione}{unbreakable, colback=orange!5, colframe=orange!70!black, boxrule=0.6pt, left=6pt, right=6pt, top=3pt, bottom=3pt, before skip=11pt plus 2pt, after skip=11pt plus 2pt}
\tcolorboxenvironment{lemma}{unbreakable, colback=orange!5, colframe=orange!70!black, boxrule=0.6pt, left=6pt, right=6pt, top=3pt, bottom=3pt, before skip=11pt plus 2pt, after skip=11pt plus 2pt}

\renewcommand{\arraystretch}{1.25}
% ---------- tipografia: niente righe orfane, titoli non isolati a fondo pagina ----------
\raggedbottom
\emergencystretch=3em
\widowpenalty=10000 \clubpenalty=10000
\pretocmd{\section}{\Needspace{10\baselineskip}}{}{}
\pretocmd{\subsection}{\Needspace{8\baselineskip}}{}{}

% ---------- macro ----------
\newcommand{\vect}[1]{\bm{#1}}
\newcommand{\dd}{\mathrm{d}}
\newcommand{\ii}{\mathrm{i}}
\newcommand{\R}{\mathbb{R}}
\newcommand{\C}{\mathbb{C}}
\newcommand{\dif}{\mathop{}\!\mathrm{d}}
\newcommand{\FT}{\mathcal{F}}
\newcommand{\ft}[1]{\tilde{#1}}
\newcommand{\abs}[1]{\left\lvert #1 \right\rvert}
\newcommand{\norm}[1]{\left\lVert #1 \right\rVert}
\newcommand{\E}{\mathbb{E}}
\DeclareMathOperator{\Var}{Var}
\DeclareMathOperator{\sinc}{sinc}
\DeclareMathOperator{\rect}{rect}
\DeclareMathOperator{\circfn}{circ}
\DeclareMathOperator{\comb}{comb}
\newcommand{\arcsec}{\mathrm{arcsec}}

% figura presa dalle slide: \slidefig[larghezza]{file}{didascalia}{etichetta}
% «file» è il nome in figure/ (estratto o ritagliato da processa_slide.py); un numero di pagina `012` indica il vecchio slide_012.jpg
\newcommand{\slidefig}[4][0.7]{%
  \begin{figure}[htbp]\centering
  \IfFileExists{figure/#2}%
    {\includegraphics[width=#1\textwidth,height=0.42\textheight,keepaspectratio]{figure/#2}}%
    {\includegraphics[width=#1\textwidth,height=0.42\textheight,keepaspectratio]{figure/slide_#2.jpg}}
  \caption{#3}\label{#4}
  \end{figure}}

% figura di una lezione: \figuralezione[larghezza]{nome-file.png}{didascalia}{etichetta}
% Se il file non è ancora in figure/, mostra un riquadro con il nome: la dispensa compila comunque.
\newcommand{\figuralezione}[4][0.8]{%
  \begin{figure}[htbp]\centering
  \IfFileExists{figure/#2}%
    {\includegraphics[width=#1\textwidth,height=0.42\textheight,keepaspectratio]{figure/#2}}%
    {\fbox{\parbox{#1\textwidth}{\centering\vspace{3em}\textit{Figura da inserire:} \texttt{\detokenize{#2}}\vspace{3em}}}}
  \caption{#3}\label{#4}
  \end{figure}}

% contenuto della dispensa senza riscontro nelle slide (comando /aderenza):
% \begin{extraslide} ... \end{extraslide}  oppure  \extraslide{testo breve}
\newtcolorbox{extraslide}{breakable, colback=red!3, colframe=red!60!black, fonttitle=\bfseries\small, title={[CONTENUTO EXTRA NON PRESENTE NELLE SLIDE]}}
